 % ============================================================
%  SECTION 8: CONCLUSIONS
% ============================================================

\section{Conclusions and Future Perspectives}

The conclusions of this review are presented in four complementary parts, covering the scientific interpretation of honey maturity, the regulatory implications of the Apimondia position, the proposed functional maturity assessment framework, and the limitations requiring further research.

This article undertook an analysis of the Apimondia position
on the production and distribution of immature honey, with
particular attention to the microbiological, regulatory, and
market implications of using osmophilic yeast count and the
morphological criterion of comb cell capping as quality indicators.
The conducted analysis leads to a number of conclusions of varying
degrees of generalization, from specific biochemical findings
to normative recommendations of global scope.

%,----------------------------------------------------------
\subsection{Main Scientific conclusions}
%,----------------------------------------------------------

\textbf{(C1) The presence of osmophilic yeasts in honey is a biologically normal phenomenon.}
Yeasts of the genera \textit{Zygosaccharomyces}, \textit{Schizosaccharomyces},
and \textit{Metschnikowia} are an inseparable component of the hive
ecosystem, regardless of production method, hive system, and country
of origin. The natural yeast population count in correct, non-fermenting
commercial honey may amount to several hundred to a few thousand CFU/g
without any effect on the microbiological stability of the product,
provided water activity $a_w < 0.60$~\cite{ruegg1981}.
This conclusion is well established in the literature for decades
and confirmed by all current systematic
reviews~\cite{biochemreactions2022,rodriguezmachado2024}.

\textbf{(C2) Water activity, not moisture content or yeast count, is the primary fermentation risk parameter.}
The physiological growth threshold of osmophilic yeasts
($a_w = 0.61$--$0.62$) defines a biologically justified
% #R1.5: "thermodynamically impossible" softened — the scientific basis is sound but absolute language is challenged by reviewer; rephrased to reflect experimental evidence
boundary below which honey fermentation is \sout{thermodynamically
impossible} \textcolor{blue}{thermodynamically unfeasible and has not been
observed under standard storage and temperature conditions,}
regardless of the yeast count expressed in
CFU/g~\cite{waaw2006,ruegg1981}.
Glucose crystallization constitutes an additional factor that
dynamically modifies the $a_w$ of the liquid phase, meaning that
fermentation risk assessment must take into account the phase state
of the product and cannot be based solely on the global water
content~\cite{zamora2006,ijfans2019}.

\textbf{(C3) Honey ripening is a continuous biological process, not a binary event.}
Maturity indicators -- invertase, glucose oxidase, and diastase
activity, organic acid profile, fructose-to-glucose ratio --
change gradually during the ripening process and do not correlate
unambiguously with the moment of cell capping by
bees~\cite{zhang2021,sun2021}.
Uncapped honey may exhibit a chemical profile appropriate for
mature honey if moisture is low; capped honey may lose maturity
characteristics as a result of inappropriate heat
treatment~\cite{codex_stan12,guo2020,zhang2021}.

\textbf{(C4) Some honey maturity criteria are derived mainly from temperate-climate, frame-hive production systems and may have limited applicability under diverse global production and microclimatic conditions.}
The climatic, technological, and biological conditions of honey
production in tropical and subtropical zones, including high
environmental humidity, log and basket hive systems, and the
distinct properties of stingless bee honeys (Meliponini), mean
that uniform morphological (capping) and microbiological (CFU/g
without $a_w$) criteria cannot serve as global quality indicators
without systematically favoring honeys produced in European
systems~\cite{tadele2025,stinglessyeastbrazil2021,sciendo2012}.
This position is consistent with the general trend in food
legislation toward recognizing the diversity of local production
systems and moving away from norms derived exclusively from
Western European practices~\cite{phiri2022}.

\textbf{(C5) The dominant form of honey adulteration on the EU market is sugar syrup adulteration, not the distribution of honey with elevated yeast count.}
Data from the coordinated control action of the European Commission
``From the Hives'' (2021--2022) unambiguously indicate that 46\%
of suspect imported honey shipments contained additions of
C$_3$/C$_4$ syrups or oligosaccharides~\cite{eufromhives2021}.
None of the available control reports identifies elevated CFU/g
count as the primary indicator of adulteration in inspection
practice. Allocation of inspection resources toward $a_w$
measurement is therefore more appropriate to the actual market
risk profile than intensification of microbiological
methods~\cite{schoder2026,eufromhives2021}.

%,----------------------------------------------------------
\subsection{Regulatory and Practical Implications of the Apimondia Position}
%,----------------------------------------------------------

The Apimondia position of 2023 on immature honey is a valuable and
necessary contribution to the protection of the integrity of the
beekeeping market. The threat posed by systemically immature honey,
produced by the method of industrial nectar dehydration in a short
time and without maintaining the biological ripening processes, is
real and constitutes unfair competition with producers adhering to
quality standards. In this respect, the Apimondia position addresses an important and legitimate concern related to honey integrity and fair trade..

This article does not undermine this intention, but indicates that
the current formulation of the position is analytically incomplete
in two important areas. First, it does not distinguish between
(a)~fermentation as a product characteristic at the time of harvest
and (b)~fermentation as a post-harvest event resulting from supply
chain errors. Second, it does not take into account the fundamental
role of water activity as a parameter superior to yeast count in the
assessment of fermentation risk. Supplementing the Apimondia position
with these two elements will not weaken its protective message, but
will instead strengthen it scientifically and make it more resistant
to legal challenges in trade disputes with importers from tropical
countries~\cite{ruegg1981,tadele2025,analytica2020}.

Thirdly, it may cause discrimination against honeys (including
European ones) for which a naturally high osmophilic yeast content
is frequent/characteristic, and which is not caused by fermentation.

%,----------------------------------------------------------
\subsection{Functional Maturity Assessment Framework}
%,----------------------------------------------------------

The proposed functional maturity assessment framework (FMAM) offers a conceptual and operational basis for further development of honey quality assessment because it:

\begin{itemize}
    \item is \textbf{scientifically justified}, based on the
          thermodynamic properties of water activity, enzymatic
          ripening kinetics, and
          metabolomics~\cite{ruegg1981,zhang2021,guo2020};
    \item is \textbf{compliant with applicable regulations}
          - Codex STAN 12-1981 and EU Directive 2001/110/EC,
          and does not require their revision, only interpretive
          extension~\cite{codex_stan12,eudir2001110};
    \item is \textbf{technically accessible}, as $a_w$ measurement
          and refractometry can be performed in any control
          laboratory without specialized
          infrastructure~\cite{waaw2006,zamora2006};
    \item is \textbf{has potential global applicability}, since criteria based
          on $a_w$ and enzymatic profile are independent of the
          hive system and climatic conditions of the place of
          production~\cite{tadele2025,stinglessyeastbrazil2021};
    \item is \textbf{proportionate to the risk}, preventing
          escalation to costly
          methods~\cite{schoder2026,nmrhoney2023}.
\end{itemize}

%,----------------------------------------------------------
\subsection{Limitations and Directions for Future Research}
%,----------------------------------------------------------

The authors of the article are aware of the limitations of this
study. The proposed FMAM model is based on empirical data available
at the time of manuscript preparation and requires further
verification through prospective multilaboratory studies.
The following directions for further research are particularly
indicated:

\begin{enumerate}
    \item \textbf{Multilaboratory validation of the decenedioic
          acid biomarker} on samples representing the global
          diversity of botanical types and production systems,
          extending beyond the three types (rapeseed, acacia,
          jujube) studied by Sun et al.~\cite{sun2021}; validation
          should encompass stingless bee honeys (Meliponini)
          and honeys from traditional East African
          systems~\cite{tadele2025};

    \item \textbf{Development and validation of a multidimensional
          functional maturity index} (FMI,
          \textit{Functional Maturity Index}) integrating the
          results of Level~II parameters into a single index
          with a value of 0--100, calibrated for different
          botanical and geographical types; such an index could
          replace binary disqualification decisions with a
          gradient approach proportionate to the deviation from
          the optimal state~\cite{zhang2021,guo2020,sun2021};

    \item \textbf{Harmonization of ISO TC~34/SC~17 standards
          taking into account the meliponiculture production
          system} -- the stingless bee honey category requires
          a separate standardization track with its own moisture
          thresholds, $a_w$, and enzymatic criteria, taking
          into account the biological properties of these
          insects~\cite{stinglessyeastbrazil2021,stinglesswjava2024}.
\end{enumerate}

Undertaking this research is possible only within the framework
of international cooperation encompassing scientific institutions,
beekeeping organizations, and regulatory bodies from countries of
both the ``Global North'' and tropical countries that are the
main honey producers. This dialogue is not only a scientific need,
but a necessary condition for building a globally equitable and
scientifically credible quality assessment system for one of the
oldest food products manufactured by humankind.
